\documentclass[11pt]{article}

\usepackage[margin=25mm]{geometry}

\usepackage{amsmath,amssymb,fontspec,unicode-math}

\IfFontExistsTF{DejaVu Serif}{\setmainfont{DejaVu Serif}}{\setmainfont{Latin Modern Roman}}

\IfFontExistsTF{Latin Modern Math}{\setmathfont{Latin Modern Math}}{\setmathfont{DejaVu Math TeX Gyre}}

\usepackage{xurl,hyperref}

\setlength{\parindent}{0pt}\setlength{\parskip}{6pt}

\begin{document}

\title{Exploring Neural Granger Causality with xLSTMs: Unveiling Temporal Dependencies in Complex Data}

\author{ELVIA preliminary AI analysis}\date{21 October 2026}\maketitle

Preliminary assessment: Minor.\par

Independent scientific review remains pending. Author adoption has not been requested.

The public weighted shrinkage is not its intended proximal map; cross-station normalization changes the graph’s interpretation. Released Lorenz sampling/noise and training settings also differ from the methods description. Historical benchmark effects remain unresolved.

\section*{Review scope}

Claim-focused reading of the original paper and relevant author-released source, with direct mathematical and static-source checks. No author program, historical verifier, model training or published benchmark was executed.

\section*{Problem 1: The released weighted shrinkage is not the proximal map}

Source: Section 3.2, Eq.(4), PDF p.5; Algorithm 1, line 16, PDF p.21; models/clstm.py, lines 342–360

Assessment: supported; impact: minor

\subsection*{Source claim (paraphrase)}

For fixed positive normalized weights, the compression step should minimize \(J(U)=\tfrac12\|U-W\|_F^2+\sum_jt_j\|U_j\|_2\). The release instead divides a column by \(\max(\|W_j\|_2,\lambda\eta)\), even when its weighted threshold is smaller.

\subsection*{Analysis}

Take two columns \(W=(1/2,2)\), \(\lambda\eta=1\), and \(\alpha=(1/10,9/10)\), embedded in the allowed hidden dimension by zero padding. Grant the release’s threshold convention \(t=2\alpha=(1/5,9/5)\). It returns \((0.15,0.20)\), whereas columnwise soft thresholding gives \((0.30,0.20)\). The objective values are respectively \(2.07125\) and \(2.06\); the first-coordinate derivative at the release output is \(0.15-0.5+0.2=-0.15\), so the return is not the unique minimizer. Uniform weights are a passing special case.

\subsection*{Finding and scope}

The public nonuniform-weight operator differs from the correct printed formula. Its later optimization and benchmark effects are unmeasured.

\subsection*{Proposed correction}

Use \((1-t/r)_+w\) for \(r=\|w\|_2>0\), with zero output for a zero column. Cauchy–Schwarz reduces the objective to \(\tfrac12(s-r)^2+ts\), uniquely minimized by \(s=(r-t)_+\). Match the threshold scaling to the stated objective.

A proposed correction is described; author adoption is not established.

\section*{Problem 2: Cross-station normalization changes the channel histories}

Source: Appendix C, Moléne description, PDF p.22; prepare\_data.py, lines 66–84

Assessment: supported; impact: minor

\subsection*{Source claim (paraphrase)}

The released Moléne loader normalizes across stations at each time: \(Y_{j,t}=(X_{j,t}-\min_\ell X_{\ell,t})/(\max_\ell X_{\ell,t}-\min_\ell X_{\ell,t})\). A graph of these channels need not describe raw-station history relations.

\subsection*{Analysis}

Use independent stationary channels \(X_t=(0,U_t,1,M_t)\), with \(U_t\) iid uniform on \([1,2]\) and \(M_t=3+M_{t-1}/2+\eta_t\), where \(\eta_t\) is iid uniform on \([-1,1]\). Its stationary solution lies in \([4,8]\). No raw channel causes another. Normalization gives \(Y_t=(0,B_t,C_t,1)\), with \(B_t=U_t/M_t\) and \(C_t=1/M_t\). The marker history reveals the previous autoregressive state, and the conditional prediction is \[\mathbb E[B_t\mid M_{t-1}=m]=\frac34\log\frac{4+m/2}{2+m/2}.\] It is strictly decreasing, taking approximately \(0.275794\) at \(m=5\) and \(0.232616\) at \(m=7\).

Even given \(M_{t-2}=m_0\) and the complete earlier target history, observing \(B_{t-1}=b\) leaves the fresh state on the nontrivial interval \((a-1,a+1)\cap(1/b,2/b)\), where \(a=3+m_0/2\). Its conditional density is proportional to \(m\). Thus the strictly varying conditional prediction is not recoverable from the target’s entire history, while the marker reveals it. This establishes a newly created predictive relation, not merely one-step correlation.

\subsection*{Finding and scope}

The normalization is legitimate for a graph of relative temperatures. Inferring raw physical station relations requires a separate justification; no actual weather edge is disproved here.

\subsection*{Proposed correction}

For raw-channel history preservation, fixed invertible channelwise affine transforms \(Y_{j,t}=a_jX_{j,t}+b_j\), \(a_j\ne0\), preserve each history sigma-field. Otherwise identify the target as the transformed-variable graph.

A proposed correction is described; author adoption is not established.

\section*{Problem 3: The public Lorenz generator uses a different observation protocol}

Source: Appendix C, Lorenz-96, PDF p.22; synthetic.py, lines 60–79; prepare\_data.py, lines 93–97

Assessment: supported; impact: minor

\subsection*{Source claim (paraphrase)}

The paper specifies output spacing \(0.05\) and noise added at each simulation step. The loader retains the generator’s default \(\delta=0.1\), and the source adds noise after deterministic integration.

\subsection*{Analysis}

With \(N=T+B\) inclusive grid points from zero to \(N\delta\), actual output spacing is \(N\delta/(N-1)\). At \(T=500\) and \(B=1000\), this is \(150/1499\approx0.100067\). The separate post-integration addition makes observational noise: unlike \(z_{k+1}=az_k+\epsilon_k\), an earlier observation perturbation cannot propagate into a later deterministic state. If “each step” meant observational noise, that portion is descriptive ambiguity rather than a stochastic-model error.

\subsection*{Finding and scope}

The inspected generator does not match the reported spacing; noise placement is explicit in the source. Which version generated the published scores remains unresolved.

\subsection*{Proposed correction}

Pin the historical generator and choose output times \(t_k=k\Delta\) for the intended spacing. Specify observational versus process noise before any experiment revision; an arbitrary stochastic-solver substitution is not justified.

A proposed correction is described; author adoption is not established.

\section*{Problem 4: The release does not honor the described batch and optimizer settings}

Source: Section 4, Training and Evaluation Details, PDF p.7; models/clstm.py, lines 497–568; F40T500.yaml

Assessment: supported; impact: minor

\subsection*{Source claim (paraphrase)}

The paper describes full-batch Adam, weight decay \(0.1\), learning-rate warmup \(2000\) and coefficient warmup \(1500\). The public path uses replacement minibatches and constructs AdamW without forwarding the declared decay.

\subsection*{Analysis}

The sampler requests \(256\,\mathrm{max\_iter}\) replacement draws and the loader hardcodes batch size \(256\), regardless of the exposed argument. With \(T=500\), context \(10\) and stride \(1\), there are \(490\) windows, so such an update is not full batch. The optimizer call forwards only learning rate; logging a decay value in extra keyword arguments does not apply it. The released configuration uses warmup \(2200\) and coefficient start \(2000\). Generic minibatch pseudocode does not resolve the explicit full-batch experimental description.

\subsection*{Finding and scope}

These are distinct reproducibility settings. Neither stochastic optimization nor AdamW is intrinsically invalid, and published-result correspondence is not established.

\subsection*{Proposed correction}

Expose and forward the selected optimizer parameters, honor an explicit full-batch or minibatch mode, and retain the resolved experiment configuration. No replacement training run is validated.

A proposed correction is described; author adoption is not established.

\section*{Limits}

Preliminary AI assessment; independent human scientific adjudication remains pending.\par

Analytic examples and release discrepancies do not remeasure the paper’s reported experiments.\par

The pinned public release is not established as the program that generated every original checkpoint or benchmark score.\par

\section*{Source identity}

arXiv:2502.09981v5 (current author version confirmed from official metadata)

\url{https://arxiv.org/pdf/2502.09981v5}

Reviewed PDF SHA-256: \nolinkurl{5163669cfdbfd24b2dbac848d2d593195d1036ac940fe257ac59bd8c9f822825}

The SHA-256 identifies the exact retained PDF inspected in this rapid source/argument check. Original supplied report bytes are retained locally. No unambiguous original PDF retrieval timestamp was recovered; file modification time was not substituted. Current official arXiv metadata checked at 2026-10-21T05:36:35.986302+00:00 identifies 2502.09981v5; implicated passages were checked in the current identified material.

\end{document}
